\chapter{Requirements Analysis and Project Specifications}\label{ch:requirements}

\section{Functional requirements}

\begin{table}[H]
\centering
\caption{Functional requirements}
\begin{tabularx}{\linewidth}{lL}
\toprule
ID & Requirement \\
\midrule
F1 & Store the billing history so that monthly and annual reports read only the period requested. \\
F2 & Apply the retention policy automatically: hot (12 months), online (36 months), archive (until 60 months), purge after. \\
F3 & Keep archived data queryable (audit, legal requests) through a single access point. \\
F4 & Allow the retention parameters to change without modifying code. \\
F5 & Provide a dry-run mode that shows what the job would do. \\
F6 & Log every lifecycle action: what, when, how many rows, how long, result. \\
F7 & Prove the integrity of billing data after every operation. \\
\bottomrule
\end{tabularx}
\end{table}

\section{Technical and non-functional requirements}

\begin{table}[H]
\centering
\caption{Technical requirements}
\begin{tabularx}{\linewidth}{lL}
\toprule
ID & Requirement \\
\midrule
T1 & Lifecycle operations must avoid copying or deleting rows one by one where possible. \\
T2 & A failure in the middle of an operation must not lose data, and a re-run must recover. \\
T3 & Data inside the legal retention can never be purged, whatever parameter is passed. \\
T4 & Old billing data cannot be modified (by accidental updates or late inserts). \\
T5 & Least privilege: the billing application cannot delete history or run lifecycle jobs, and analysts are read-only. \\
T6 & No credentials in source code or in the Git repository. \\
T7 & The whole solution must be reproducible from scripts on the course Docker environment. \\
T8 & The lab must not exhaust the laptop disk. \\
T9 & Use only licensed features, and document the licence impact. \\
\bottomrule
\end{tabularx}
\end{table}

\section{Retention policy}

The legal and business rules translate into four tiers, stored as parameters in the
\code{DLM\_CONFIG} table rather than hard-coded:

\begin{table}[H]
\centering
\caption{Retention policy parameters}
\begin{tabularx}{\linewidth}{llL}
\toprule
Parameter & Value & Meaning \\
\midrule
\code{HOT\_MONTHS} & 12 & months kept uncompressed and writable \\
\code{ONLINE\_MONTHS} & 36 & months kept in the online billing table (hot + warm) \\
\code{RETENTION\_MONTHS} & 60 & legal retention; older months are purged \\
\code{MIN\_RETENTION\_MONTHS} & 60 & safety floor: no purge may use a shorter retention \\
\code{AS\_OF\_MONTH} & 2026-11 & reference ``today'', so that runs are reproducible \\
\bottomrule
\end{tabularx}
\end{table}

\section{Measurable indicators and acceptance criteria}

We formalised the indicators after a first exploratory benchmark and before the final run
whose numbers this report presents. The brief asks for indicators to be defined before
testing, so we state this order openly. The indicators compare the two designs on the same
data, with the same SQL text and bind values.

\begin{table}[H]
\centering
\caption{Key performance indicators and targets}
\label{tab:kpi}
\begin{tabularx}{\linewidth}{lXl}
\toprule
KPI & Indicator & Target \\
\midrule
K1 & Month-end report (Q1): elapsed time and logical reads & $\geq$5$\times$ time, $\geq$10$\times$ LIO \\
K2 & Audit query on an archived month (Q7) & $\geq$5$\times$ fewer LIO \\
K3 & OLTP lookups (Q2 by customer, Q5 by id) & no time regression $>$2$\times$ \\
K4 & Archive one month (and the whole monthly job) & $\geq$10$\times$ faster, $<$5~s \\
K5 & Purge one month & $<$1~s \\
K6 & Monthly load + statistics refresh & not slower than baseline \\
K7 & Storage of the retained data & $\geq$25\,\% smaller \\
K8 & Data integrity vs load-time control totals & 100\,\% of months \\
K9 & Recovery after a crash & 0 rows lost, recovery = re-run \\
K10 & Security controls (asserted checks) & 100\,\% \\
K11 & Automation & one command per month \\
\bottomrule
\end{tabularx}
\end{table}

The acceptance criteria of the brief are: (1) the design preserves data integrity,
(2) retention operations work correctly, (3) query or maintenance improvements are measured
on comparable workloads, and (4) licensing restrictions and technical trade-offs are
documented. Chapter~\ref{ch:performance} checks each of them.
